\documentclass[a4paper,12pt]{paper}

\usepackage{mycommands}

% Worked-example layout: boxed problem statement followed by a solution.
\newcounter{example}
\renewenvironment{example}[1]{%  % the paper class already defines "example"
  \refstepcounter{example}%
  \begin{mdframed}[linewidth=1pt,innertopmargin=6pt,innerbottommargin=6pt]%
  \noindent\textbf{Example \theexample: #1.}\ }{%
  \end{mdframed}}
\newenvironment{solution}{\par\noindent\textbf{Solution.}\ }{\par\medskip}

\title{Worked examples for Lecture 21: Equilibrium figures}
\author{David Al-Attar \\ Michaelmas Term, \the\year}

\begin{document}

\maketitle

\noindent These examples accompany Lecture~21. They are intended as practice in the concrete
calculations that underlie the lecture, and are less demanding than the questions on the
problem sets. Numerical results were checked with the scripts in the \texttt{python}
directory. Throughout, $G=6.674\times10^{-11}\,\mathrm{m^{3}\,kg^{-1}\,s^{-2}}$, the mean radius
of the Earth is $b=6371\,\mathrm{km}$, its mass is $M=5.972\times10^{24}\,\mathrm{kg}$, its
mean density is $\overline{\rho}=5515\,\mathrm{kg\,m^{-3}}$, and its angular velocity is
$\Omega=7.292\times10^{-5}\,\mathrm{s^{-1}}$. Values quoted for the Preliminary Reference Earth
Model (PREM) are from Dziewonski \& Anderson (1981).

\begin{example}{Pressure inside a homogeneous non-rotating planet}
A non-rotating planet of radius $b$ has uniform density $\rho$. Show that the hydrostatic
pressure is
\begin{equation*}
p(r)=\frac{2}{3}\pi G\rho^{2}(b^{2}-r^{2}),
\end{equation*}
and express the central pressure in terms of the mass $M$ and the surface gravity $g(b)$.
Evaluate the central pressure for a planet with the mean radius and mean density of the Earth,
compare with the value $3.64\times10^{11}\,\mathrm{Pa}$ at the centre of PREM, and explain the
discrepancy.
\end{example}

\begin{solution}
For a spherically symmetric non-rotating planet the hydrostatic equilibrium equation reduces to
\begin{equation}
\frac{\ddns p}{\ddns r}+\rho\,g_{0}=0,\quad g_{0}(r)=\frac{\ddns\phi_{0}}{\ddns r}=\frac{GM(r)}{r^{2}},
\label{eq:1}
\end{equation}
where $M(r)$ is the mass within radius $r$, and $g_{0}>0$ is the magnitude of the gravitational
acceleration. The expression for $g_{0}$ follows from Gauss's law applied to Poisson's equation, or
equivalently from Newton's shell theorem. For uniform density $M(r)=\frac{4}{3}\pi\rho r^{3}$, and so
\begin{equation}
g_{0}(r)=\frac{4}{3}\pi G\rho\,r,
\label{eq:2}
\end{equation}
which increases linearly from the centre to the surface. Integrating eq.~(\ref{eq:1}) inwards from
the surface, where $p(b)=0$,
\begin{equation}
p(r)=\int_{r}^{b}\rho\,g_{0}(r')\,\dd r'=\frac{4}{3}\pi G\rho^{2}\int_{r}^{b}r'\,\dd r'=\frac{2}{3}\pi G\rho^{2}(b^{2}-r^{2}),
\label{eq:3}
\end{equation}
as required. The pressure is a parabola in $r$, and its central value is
\begin{equation}
p(0)=\frac{2}{3}\pi G\rho^{2}b^{2}.
\label{eq:4}
\end{equation}
Substituting $\rho=3M/(4\pi b^{3})$ we can write this as
\begin{equation}
p(0)=\frac{3GM^{2}}{8\pi b^{4}}=\frac{3\,g(b)^{2}}{8\pi G},
\label{eq:5}
\end{equation}
where $g(b)=GM/b^{2}$ is the surface gravity. The second form is worth remembering: it says that
the central pressure of a homogeneous planet is fixed by its surface gravity alone, and does not
otherwise depend on its size.

For the Earth's mean density and radius, eq.~(\ref{eq:4}) gives
\begin{equation}
p(0)=\frac{2}{3}\pi\times6.674\times10^{-11}\times(5515)^{2}\times(6.371\times10^{6})^{2}\,\mathrm{Pa}
=1.73\times10^{11}\,\mathrm{Pa},
\label{eq:6}
\end{equation}
which is less than half the PREM value of $3.64\times10^{11}\,\mathrm{Pa}$. The discrepancy arises
because the Earth's mass is concentrated towards its centre: the density in the core is roughly
twice that in the mantle, and within each region density increases with depth through
self-compression. Two effects then act in the same direction. First, concentrating mass towards
the centre increases $M(r)/r^{2}$ at depth, so that $g_{0}$ within the Earth is everywhere larger
than the linear profile of eq.~(\ref{eq:2}); in PREM, $g_{0}$ is roughly constant at about
$10\,\mathrm{m\,s^{-2}}$ through the whole mantle, and only falls towards zero within the core.
Second, the weight per unit volume $\rho g_{0}$ at depth is further increased by the larger
density there. Both effects steepen the pressure gradient at depth, and hence raise the central
pressure. The homogeneous model gets the mass and surface gravity right, but the pressure
within a planet depends on how that mass is distributed with radius, and a homogeneous model
places too much of it near the surface. The next example quantifies this with the simplest
model that has a dense core.
\end{solution}

\begin{example}{Two-layer planet}
A non-rotating planet consists of a core of uniform density $\rho_{\mathrm{c}}$ and radius
$c$, surrounded by a mantle of uniform density $\rho_{\mathrm{m}}$ extending to the surface
$r=b$. Obtain the gravitational acceleration $g_{0}(r)$ and the hydrostatic pressure $p(r)$ in
both layers. Taking $\rho_{\mathrm{c}}=11\times10^{3}\,\mathrm{kg\,m^{-3}}$,
$\rho_{\mathrm{m}}=4.5\times10^{3}\,\mathrm{kg\,m^{-3}}$, $c=3480\,\mathrm{km}$, and
$b=6371\,\mathrm{km}$, compute the mass, the moment of inertia factor $I/(Mb^{2})$, the pressure
at the core--mantle boundary (CMB), and the central pressure. Compare with the Earth's values
$I/(Mb^{2})=0.3307$, and with PREM's $p(\mathrm{CMB})=1.36\times10^{11}\,\mathrm{Pa}$ and
$p(0)=3.64\times10^{11}\,\mathrm{Pa}$.
\end{example}

\begin{solution}
The mass within radius $r$ is
\begin{equation}
M(r)=\begin{cases}\frac{4}{3}\pi\rho_{\mathrm{c}}r^{3} & 0\le r\le c,\\[4pt]
\frac{4}{3}\pi\left[\rho_{\mathrm{c}}c^{3}+\rho_{\mathrm{m}}(r^{3}-c^{3})\right] & c\le r\le b,\end{cases}
\label{eq:7}
\end{equation}
and so, with $g_{0}=GM(r)/r^{2}$ as in Example~1,
\begin{equation}
g_{0}(r)=\begin{cases}\frac{4}{3}\pi G\rho_{\mathrm{c}}\,r & 0\le r\le c,\\[4pt]
\frac{4}{3}\pi G\left[\rho_{\mathrm{m}}\,r+(\rho_{\mathrm{c}}-\rho_{\mathrm{m}})\dfrac{c^{3}}{r^{2}}\right] & c\le r\le b.\end{cases}
\label{eq:8}
\end{equation}
In the mantle, $g_{0}$ is the sum of the linearly increasing field of a homogeneous planet of density
$\rho_{\mathrm{m}}$ and the inverse-square field of the excess core mass
$\frac{4}{3}\pi(\rho_{\mathrm{c}}-\rho_{\mathrm{m}})c^{3}$. These two terms vary in opposite senses with $r$,
which is why $g_{0}$ is so nearly uniform through the Earth's mantle. Note that $g_{0}$ is
continuous at $r=c$, as it must be since $M(r)$ is continuous.

The pressure follows from $\ddns p/\ddns r=-\rho g_{0}$ with $p(b)=0$. In the mantle,
\begin{equation}
p(r)=\rho_{\mathrm{m}}\int_{r}^{b}g_{0}(r')\,\dd r'
=\frac{2}{3}\pi G\rho_{\mathrm{m}}^{2}(b^{2}-r^{2})+\frac{4}{3}\pi G\rho_{\mathrm{m}}(\rho_{\mathrm{c}}-\rho_{\mathrm{m}})\,c^{3}\left(\frac{1}{r}-\frac{1}{b}\right),
\quad c\le r\le b.
\label{eq:9}
\end{equation}
The first term is the pressure in a homogeneous planet of density $\rho_{\mathrm{m}}$, while the
second is the additional pressure needed to support the mantle against the attraction of the excess
core mass. Setting $r=c$ gives the pressure at the CMB,
\begin{equation}
p(c)=\frac{2}{3}\pi G\rho_{\mathrm{m}}^{2}(b^{2}-c^{2})+\frac{4}{3}\pi G\rho_{\mathrm{m}}(\rho_{\mathrm{c}}-\rho_{\mathrm{m}})\,c^{2}\left(1-\frac{c}{b}\right).
\label{eq:10}
\end{equation}
Within the core the density is uniform and $g_{0}$ is linear in $r$, so the calculation of
Example~1 applies with the constant of integration fixed by continuity of $p$ at the CMB,
\begin{equation}
p(r)=p(c)+\frac{2}{3}\pi G\rho_{\mathrm{c}}^{2}(c^{2}-r^{2}),\quad 0\le r\le c.
\label{eq:11}
\end{equation}
The pressure is continuous at $r=c$, but its gradient $-\rho g_{0}$ jumps by the factor
$\rho_{\mathrm{c}}/\rho_{\mathrm{m}}$, since $g_{0}$ is continuous while $\rho$ is not.

We now put in the numbers. From eq.~(\ref{eq:7}) with $r=b$, the mass is
\begin{equation}
M=\frac{4}{3}\pi\left[\rho_{\mathrm{c}}c^{3}+\rho_{\mathrm{m}}(b^{3}-c^{3})\right]=6.02\times10^{24}\,\mathrm{kg},
\label{eq:12}
\end{equation}
within one per cent of the Earth's mass. The moment of inertia about any axis through the centre of
a spherically symmetric body is $I=\frac{8}{3}\pi\int_{0}^{b}\rho\,r^{4}\,\dd r$, and hence
\begin{equation}
I=\frac{8}{15}\pi\left[\rho_{\mathrm{c}}c^{5}+\rho_{\mathrm{m}}(b^{5}-c^{5})\right]=8.47\times10^{37}\,\mathrm{kg\,m^{2}},
\quad\frac{I}{Mb^{2}}=0.347.
\label{eq:13}
\end{equation}
For a homogeneous sphere $I/(Mb^{2})=0.4$, and the observed value of $0.3307$ for the Earth is
one of the classical indications that its density increases with depth. The two-layer model
accounts for most of the reduction, but not all of it, because within each layer the real
density increases with depth under compression, which the model neglects. The gravitational
acceleration from eq.~(\ref{eq:8}) is $g_{0}(b)=9.90\,\mathrm{m\,s^{-2}}$ at the surface and
$g_{0}(c)=10.70\,\mathrm{m\,s^{-2}}$ at the CMB, the values in PREM being $9.82$ and
$10.68\,\mathrm{m\,s^{-2}}$. For the pressure, the two terms in eq.~(\ref{eq:10}) evaluate to
$8.06\times10^{10}\,\mathrm{Pa}$ and $4.49\times10^{10}\,\mathrm{Pa}$, giving
\begin{equation}
p(c)=1.26\times10^{11}\,\mathrm{Pa},
\label{eq:14}
\end{equation}
and the core then adds $\frac{2}{3}\pi G\rho_{\mathrm{c}}^{2}c^{2}=2.05\times10^{11}\,\mathrm{Pa}$,
so that from eq.~(\ref{eq:11})
\begin{equation}
p(0)=3.30\times10^{11}\,\mathrm{Pa}.
\label{eq:15}
\end{equation}
These are within $8$ and $9$ per cent of the PREM values of $1.36\times10^{11}\,\mathrm{Pa}$ and
$3.64\times10^{11}\,\mathrm{Pa}$, and almost twice the central pressure of the homogeneous model
of Example~1. The model errs on the low side for the same reason that $I/(Mb^{2})$ errs on the
high side: neglecting self-compression places too little mass at depth. A model with only two free
densities nevertheless reproduces the mass, moment of inertia, and pressures of the Earth to within
about ten per cent, which illustrates how strongly the gross structure of a planet is controlled by
the existence of a dense core.
\end{solution}

\begin{example}{Centrifugal potential and $Y_{20}$}
For $\bm{\Omega}=\Omega\hat{\mathbf{e}}_{3}$ show that the centrifugal potential
$\psi=\frac{1}{2}(\Omega_{i}\Omega_{j}-\Omega_{k}\Omega_{k}\delta_{ij})x_{i}x_{j}$ defined in the
lecture can be written
\begin{equation*}
\psi=-\frac{1}{3}\Omega^{2}r^{2}+\frac{1}{3}\Omega^{2}r^{2}P_{2}(\cos\theta),
\end{equation*}
with $\theta$ the colatitude and $P_{2}$ the Legendre polynomial of degree two, and hence that its
aspherical part is proportional to $r^{2}Y_{20}$. Explain why the spherical part can be absorbed
into the reference model, and deduce that all first-order perturbations of a slowly rotating
hydrostatic planet have $Y_{20}$ angular dependence. Using the lecture's definition of the
ellipticity, show that the surface of the planet is
$r(\theta)=b\left[1-\frac{2}{3}\epsilon(b)P_{2}(\cos\theta)\right]$ and verify that
$\epsilon(b)$ is the difference between the equatorial and polar radii divided by the mean radius.
\end{example}

\begin{solution}
With $\Omega_{i}=\Omega\,\delta_{i3}$ the definition gives
\begin{equation}
\psi=\frac{1}{2}\left(\Omega^{2}x_{3}^{2}-\Omega^{2}x_{k}x_{k}\right)=-\frac{1}{2}\Omega^{2}\left(x_{1}^{2}+x_{2}^{2}\right)
=-\frac{1}{2}\Omega^{2}r^{2}\sin^{2}\theta,
\label{eq:16}
\end{equation}
which is minus one half of the square of the distance from the rotation axis times $\Omega^{2}$.
Its gradient is $-\Omega^{2}(x_{1},x_{2},0)$, the outward centrifugal acceleration
$-\bm{\Omega}\times(\bm{\Omega}\times\mathbf{x})$, as it should be. Recalling that
$P_{2}(\cos\theta)=\frac{1}{2}(3\cos^{2}\theta-1)$, we have
\begin{equation}
\sin^{2}\theta=1-\cos^{2}\theta=1-\frac{1}{3}\left[2P_{2}(\cos\theta)+1\right]=\frac{2}{3}\left[1-P_{2}(\cos\theta)\right],
\label{eq:17}
\end{equation}
and substituting into eq.~(\ref{eq:16}) gives
\begin{equation}
\psi=-\frac{1}{3}\Omega^{2}r^{2}+\frac{1}{3}\Omega^{2}r^{2}P_{2}(\cos\theta),
\label{eq:18}
\end{equation}
as required. The zonal spherical harmonic of degree two is
$Y_{20}=\sqrt{5/(4\pi)}\,P_{2}(\cos\theta)$, with the normalisation
$\int|Y_{20}|^{2}\,\dd S=1$ over the unit sphere, and so the aspherical part of $\psi$ is
\begin{equation}
\psi^{\mathrm{asph}}=\sqrt{\frac{4\pi}{5}}\,\frac{1}{3}\Omega^{2}r^{2}\,Y_{20}.
\label{eq:19}
\end{equation}
As a check, eq.~(\ref{eq:16}) gives $\nabla^{2}\psi=-2\Omega^{2}$, and since $\nabla^{2}r^{2}=6$ the
spherical part of eq.~(\ref{eq:18}) supplies the whole of this, while $r^{2}Y_{20}$ is a solid
harmonic and hence contributes nothing.

The spherical part $-\frac{1}{3}\Omega^{2}r^{2}$ depends on $r$ alone. Adding it to the
gravitational potential of the reference model simply changes the zeroth-order equilibrium equation
to $\ddns p_{0}/\ddns r+\rho_{0}(g_{0}-\frac{2}{3}\Omega^{2}r)=0$, that is, the spherical model feels a
slightly reduced effective gravity. The relative change, $\frac{2}{3}\Omega^{2}r/g_{0}$, is at most
$\frac{2}{3}\Omega^{2}b^{3}/(GM)\approx2\times10^{-3}$ for the Earth, so this is a first-order correction
to a spherically symmetric model. It changes the reference values of $p_{0}$ and $g_{0}$ by a fraction
of order the small parameter, but does not alter the spherical symmetry of the reference state, and
it is this spherical symmetry alone that the first-order theory relies upon. We may therefore
regard the spherical part as already included in $\rho_{0}$, $p_{0}$, and $\phi_{0}$, and take
$\psi_{1}=\psi^{\mathrm{asph}}$. This is precisely the freedom used in the lecture when the
perturbations were assumed to average to zero over spherical surfaces.

Now consider the first-order problem. The boundary value problem for $\phi_{1}$ is linear in
$\phi_{1}$ and $\psi_{1}$, its coefficients $4\pi G\rho_{0}'/g_{0}$ and $4\pi G\rho_{0}/g_{0}$ depend on
$r$ alone, and the Laplacian commutes with rotations. Substituting the separated form
$\phi_{1}=f(r)Y_{20}$, the angular dependence factors out of every term, leaving an ordinary
differential equation in $r$ for $f$ whose forcing is proportional to $r^{2}$. Any other spherical
harmonic component of $\phi_{1}$ would satisfy the same problem with no forcing, and must vanish
for the perturbation to be unique.\footnote{The one exception is the degree-one part, which
corresponds to a rigid translation of the whole planet and is removed by keeping the centre of
mass at the origin.} Hence $\phi_{1}$, and therefore
$\gamma_{1}=\phi_{1}+\psi_{1}$, is proportional to $Y_{20}$. The remaining first-order fields are given
in terms of $\gamma_{1}$ by the algebraic relations $p_{1}=-\rho_{0}\gamma_{1}$, $h_{1}=-\gamma_{1}/g_{0}$,
and $\rho_{1}=(\rho_{0}'/g_{0})\gamma_{1}$, with coefficients depending on $r$ alone, and so all have the
same $Y_{20}$ dependence. Since $Y_{20}$ is independent of longitude and even in $\cos\theta$, the
perturbed planet is axisymmetric and symmetric about its equatorial plane.

Finally, the lecture defines the ellipticity by
$\gamma_{1}=\sqrt{4\pi/5}\,\frac{2}{3}rg_{0}\,\epsilon\,Y_{20}=\frac{2}{3}rg_{0}\,\epsilon\,P_{2}(\cos\theta)$,
and therefore the surface displacement is
\begin{equation}
h_{1}=-\frac{\gamma_{1}}{g_{0}}\bigg|_{r=b}=-\frac{2}{3}b\,\epsilon(b)P_{2}(\cos\theta),
\quad r(\theta)=b\left[1-\frac{2}{3}\epsilon(b)P_{2}(\cos\theta)\right].
\label{eq:20}
\end{equation}
At the equator $P_{2}(0)=-\frac{1}{2}$ and at the poles $P_{2}(\pm1)=1$, so the equatorial and polar
radii are
\begin{equation}
r_{\mathrm{eq}}=b\left[1+\frac{1}{3}\epsilon(b)\right],\quad r_{\mathrm{pol}}=b\left[1-\frac{2}{3}\epsilon(b)\right],
\label{eq:21}
\end{equation}
and their difference is $r_{\mathrm{eq}}-r_{\mathrm{pol}}=b\,\epsilon(b)$. For $\epsilon>0$ the planet is
oblate, with the equatorial bulge twice as far from the mean sphere as the polar depression, as
required by the vanishing of the spherical average of $P_{2}$. That average also shows that $b$ is the
mean radius of the perturbed surface, and that its volume is unchanged to first order. Hence
$\epsilon(b)=(r_{\mathrm{eq}}-r_{\mathrm{pol}})/b$, which is the geometrical interpretation quoted in
the lecture; for the Earth $r_{\mathrm{eq}}-r_{\mathrm{pol}}\approx21.4\,\mathrm{km}$ and
$\epsilon(b)\approx1/298$. The lesson is that the entire aspherical structure of a slowly rotating
hydrostatic planet is described by one radial function, because the forcing has a single angular
component and the problem is linear.
\end{solution}

\begin{example}{Clairaut's equation for a homogeneous planet}
Show that for a planet of uniform density $\rho_{0}$ one has $8\pi G\rho_{0}/g_{0}=6/r$, so that
Clairaut's equation reduces to $\epsilon''+6\epsilon'/r=0$, whose regular solution is
$\epsilon=$ const. Apply the surface boundary condition to obtain Newton's result
\begin{equation*}
\epsilon=\frac{5}{4}\frac{\Omega^{2}b^{3}}{GM}=\frac{5}{4}m,
\end{equation*}
where $m=\Omega^{2}b^{3}/(GM)$. Evaluate $1/\epsilon$ for a homogeneous Earth, compare with the
observed value of $298$, and explain the sign of the discrepancy.
\end{example}

\begin{solution}
For constant $\rho_{0}$ we found in Example~1 that $g_{0}=\frac{4}{3}\pi G\rho_{0}r$, and so
\begin{equation}
\frac{8\pi G\rho_{0}}{g_{0}}=\frac{8\pi G\rho_{0}}{\frac{4}{3}\pi G\rho_{0}r}=\frac{6}{r}.
\label{eq:22}
\end{equation}
More generally, $g_{0}=\frac{4}{3}\pi G\overline{\rho}(r)\,r$ with $\overline{\rho}(r)$ the mean density
inside radius $r$, and the coefficient in Clairaut's equation is $6\rho_{0}(r)/[r\overline{\rho}(r)]$;
it is the ratio of the local density to the mean density within $r$ that controls the solution.
Substituting eq.~(\ref{eq:22}) into Clairaut's equation,
\begin{equation}
\frac{\ddns^{2}\epsilon}{\ddns r^{2}}+\frac{6}{r}\left(\frac{\ddns\epsilon}{\ddns r}+\frac{\epsilon}{r}\right)-\frac{6\epsilon}{r^{2}}
=\frac{\ddns^{2}\epsilon}{\ddns r^{2}}+\frac{6}{r}\frac{\ddns\epsilon}{\ddns r}=0,
\label{eq:23}
\end{equation}
the terms in $\epsilon/r^{2}$ cancelling exactly. This is an equation of Euler type, and trying
$\epsilon=r^{n}$ gives $n(n-1)+6n=n(n+5)=0$, so the general solution is
\begin{equation}
\epsilon=A+Br^{-5}.
\label{eq:24}
\end{equation}
The second solution is singular at the centre and violates the condition
$\epsilon'(0)=0$, so we must take $B=0$: within a homogeneous planet the ellipticity is the same
on every level surface, and all the level surfaces are similar spheroids.

The constant $A$ is fixed by the surface condition. With $\epsilon'(b)=0$ the boundary condition
from the lecture,
\begin{equation}
\frac{\ddns\epsilon}{\ddns r}\bigg|_{r=b}=\frac{1}{b}\left[\frac{5\Omega^{2}b^{3}}{2GM}-2\epsilon(b)\right],
\label{eq:25}
\end{equation}
becomes $2\epsilon=5\Omega^{2}b^{3}/(2GM)$, and so
\begin{equation}
\epsilon=\frac{5}{4}\frac{\Omega^{2}b^{3}}{GM}=\frac{5}{4}m,\quad m=\frac{\Omega^{2}b^{3}}{GM}=\frac{\Omega^{2}b}{g(b)}=\frac{3\Omega^{2}}{4\pi G\rho_{0}},
\label{eq:26}
\end{equation}
which is Newton's result. The parameter $m$ is the ratio of the centrifugal to the gravitational
acceleration at the equator; the last form shows it to be, up to a numerical factor, the ratio of
centrifugal to gravitational potentials used in the lecture to define slow rotation. For the
Earth, $m=(7.292\times10^{-5})^{2}(6.371\times10^{6})^{3}/(3.986\times10^{14})=3.45\times10^{-3}$,
so that
\begin{equation}
\epsilon=\frac{5}{4}m=4.31\times10^{-3},\quad\frac{1}{\epsilon}=232,
\label{eq:27}
\end{equation}
compared with the observed $1/\epsilon=298$. The homogeneous planet is thus too flattened by
about $30$ per cent.\footnote{Newton obtained $1/230$ in the ``Principia'' using a slightly different
value of $m$. The French geodetic expeditions to Lapland and Peru in the 1730s confirmed that the
Earth is oblate, against the prolate Earth favoured by the Cassinis, but accurate values of the
flattening came only later from geodesy and from the analysis of the Moon's motion. Clairaut's
theory of 1743 explained why the Earth is less flattened than a homogeneous planet.}

To understand the sign of the discrepancy it helps to see where the factor $5/4$ comes from.
Write $\gamma_{1}=\frac{2}{3}rg_{0}\,\epsilon\,P_{2}(\cos\theta)$ as in Example~3. Displacing the surface of
a homogeneous planet by $h_{1}=-\frac{2}{3}b\,\epsilon\,P_{2}$ is equivalent, to first order, to placing
a surface mass density $\sigma=\rho_{0}h_{1}$ on the sphere $r=b$. The interior potential of a surface
density $\sigma_{2}P_{2}(\cos\theta)$ on a sphere of radius $b$ is
$-\frac{4}{5}\pi Gb\,\sigma_{2}\,(r/b)^{2}P_{2}(\cos\theta)$, a standard consequence of the Legendre
expansion of $1/\|\mathbf{x}-\mathbf{x}'\|$, and so
\begin{equation}
\phi_{1}=\frac{8}{15}\pi G\rho_{0}\,\epsilon\,r^{2}P_{2}(\cos\theta),\quad r<b.
\label{eq:28}
\end{equation}
This is positive at the poles and negative at the equator: the extra mass in the equatorial bulge
deepens the potential well there. Adding the aspherical centrifugal potential of eq.~(\ref{eq:18})
and equating to $\gamma_{1}=\frac{2}{3}rg_{0}\,\epsilon\,P_{2}=\frac{8}{9}\pi G\rho_{0}\,\epsilon\,r^{2}P_{2}$
gives
\begin{equation}
\frac{8}{9}\pi G\rho_{0}\,\epsilon=\frac{8}{15}\pi G\rho_{0}\,\epsilon+\frac{1}{3}\Omega^{2}
\quad\Longrightarrow\quad
\epsilon=\frac{15\,\Omega^{2}}{16\pi G\rho_{0}}=\frac{5}{4}m,
\label{eq:29}
\end{equation}
in agreement with eq.~(\ref{eq:26}), and providing an independent check on the signs in
Clairaut's equation and its surface condition. The proportionality to $r^{2}$ of every term is
the reason $\epsilon$ is constant. Had we ignored the self-gravity of the bulge, dropping the
first term on the right of eq.~(\ref{eq:29}), we would have found
$\epsilon=3\Omega^{2}/(8\pi G\rho_{0})=\frac{1}{2}m$, the value for a massless fluid envelope
surrounding a point mass, for which $1/\epsilon=580$ for the Earth. The self-attraction of the
equatorial bulge therefore amplifies the flattening of a homogeneous planet by a factor of
$5/2$.

The real Earth lies between these two limits. In a centrally condensed planet the equatorial
bulge contains a smaller fraction of the total mass, its self-attraction is weaker, and the
flattening is correspondingly smaller than $\frac{5}{4}m$, though larger than $\frac{1}{2}m$. The
observed $1/\epsilon=298>232$ is therefore a further indication, independent of the moment of
inertia, that the Earth's density increases with depth. Solving Clairaut's equation by shooting for
the two-layer model of Example~2 gives $1/\epsilon(b)=285$ at the surface and $1/\epsilon=374$ at
the CMB, so that $\epsilon$ increases outwards from the centre. The same trend is seen in the
lower panel of the figure in Lecture~21, where $1/\epsilon$ is about $300$ at the surface and
increases inwards, and it is a general property of solutions of Clairaut's equation for a
density that does not increase outwards. The level surfaces deep within a planet are less
flattened than its outer surface because they feel the centrifugal potential, which grows as
$r^{2}$, less strongly relative to the gravitational attraction of the mass within them.
\end{solution}

\begin{example}{Hydrostatic level surfaces}
In the lecture it was shown that in a hydrostatic planet the level surfaces of $p$, $\rho$, and
$\gamma$ coincide. Verify this directly for the first-order solution of the slowly rotating
problem, $p_{1}=-\rho_{0}\gamma_{1}$ and $\rho_{1}=(\rho_{0}'/g_{0})\gamma_{1}$, by showing that the level
surfaces of all three fields through the reference radius $r$ are displaced radially by the same
amount $-\gamma_{1}/g_{0}$. Show also that the radial component of the first-order equilibrium
equation, $\nabla p_{1}+\rho_{0}\nabla\gamma_{1}+\rho_{1}g_{0}\hat{\mathbf{r}}=\mathbf{0}$, is then
satisfied identically, and interpret the function $\epsilon(r)$ in the light of this result.
\end{example}

\begin{solution}
Consider a field $F=F_{0}(r)+s\,F_{1}(r,\theta,\varphi)+\cdots$ whose zeroth-order part depends on $r$
alone. The level surface of $F$ that passes through radius $r$ in the reference model is, to first
order, $r+s\,\delta r_{F}(\theta,\varphi)$, where
\begin{equation}
F_{0}(r+s\,\delta r_{F})+s\,F_{1}(r,\theta,\varphi)=F_{0}(r)
\quad\Longrightarrow\quad
\delta r_{F}=-\frac{F_{1}}{F_{0}'},
\label{eq:30}
\end{equation}
provided $F_{0}'\neq0$. This is the same expansion that produced the boundary condition
$p_{1}+p_{0}'h_{1}=0$ on the perturbed surface in the lecture, the outer surface being the level
surface $p=0$. Applying eq.~(\ref{eq:30}) in turn to the three fields, and using the zeroth-order
hydrostatic equation $p_{0}'=-\rho_{0}g_{0}$ together with $\gamma_{0}'=\phi_{0}'=g_{0}$, we find
\begin{align}
\delta r_{p}&=-\frac{p_{1}}{p_{0}'}=-\frac{-\rho_{0}\gamma_{1}}{-\rho_{0}g_{0}}=-\frac{\gamma_{1}}{g_{0}},\label{eq:31}\\
\delta r_{\rho}&=-\frac{\rho_{1}}{\rho_{0}'}=-\frac{1}{\rho_{0}'}\frac{\rho_{0}'}{g_{0}}\gamma_{1}=-\frac{\gamma_{1}}{g_{0}},\label{eq:32}\\
\delta r_{\gamma}&=-\frac{\gamma_{1}}{\gamma_{0}'}=-\frac{\gamma_{1}}{g_{0}}.\label{eq:33}
\end{align}
The three displacements are identical, and equal to the surface displacement $h_{1}$ of the lecture
evaluated at radius $r$. Thus at every depth the perturbed level surfaces of pressure, density, and
gravity potential coincide, as the exact theory requires. Where $\rho_{0}'=0$ the density has no
level surfaces to speak of, but there $\rho_{1}=0$ as well, and there is nothing to check.

The relations for $p_{1}$ and $\rho_{1}$ can be read in the same way. Writing
$h(r,\theta,\varphi)=-\gamma_{1}/g_{0}$ for the displacement of the level surface through $r$, we have
\begin{equation}
p_{1}=-\rho_{0}\gamma_{1}=-p_{0}'h,\quad \rho_{1}=\frac{\rho_{0}'}{g_{0}}\gamma_{1}=-\rho_{0}'h,
\label{eq:34}
\end{equation}
which are exactly the changes at a fixed point produced by displacing the spherical profiles
$p_{0}(r)$ and $\rho_{0}(r)$ radially outwards through $h$. The first-order solution is, in other
words, nothing more than the reference model with each of its level surfaces carried onto the
corresponding level surface of the perturbed gravity potential.

Now take the radial component of the first-order equilibrium equation. Substituting
$p_{1}=-\rho_{0}\gamma_{1}$,
\begin{equation}
\frac{\partial p_{1}}{\partial r}+\rho_{0}\frac{\partial\gamma_{1}}{\partial r}+\rho_{1}g_{0}
=-\rho_{0}'\gamma_{1}-\rho_{0}\frac{\partial\gamma_{1}}{\partial r}+\rho_{0}\frac{\partial\gamma_{1}}{\partial r}+\rho_{1}g_{0}
=g_{0}\left(\rho_{1}-\frac{\rho_{0}'}{g_{0}}\gamma_{1}\right),
\label{eq:35}
\end{equation}
which vanishes by the expression for $\rho_{1}$. The tangential components are satisfied because
$p_{1}+\rho_{0}\gamma_{1}=0$ identically. Note that the argument can be run in the other direction:
once the tangential components have given $p_{1}=-\rho_{0}\gamma_{1}$, the radial component yields
$\rho_{1}=(\rho_{0}'/g_{0})\gamma_{1}$ directly, without taking the curl of the equation as was done in
the lecture. The two routes must agree, because the curl of the equilibrium equation is a
consequence of the equation itself.

Finally, with $\gamma_{1}=\frac{2}{3}rg_{0}\,\epsilon(r)P_{2}(\cos\theta)$ the common displacement is
\begin{equation}
h=-\frac{\gamma_{1}}{g_{0}}=-\frac{2}{3}r\,\epsilon(r)P_{2}(\cos\theta),
\label{eq:36}
\end{equation}
which is eq.~(\ref{eq:20}) of Example~3 with $b$ replaced by $r$. Every level surface of the
perturbed planet is therefore an oblate spheroid of mean radius $r$ whose equatorial and polar
radii differ by $r\,\epsilon(r)$, and $\epsilon(r)$ is the ellipticity of the level surface of mean radius
$r$. The lecture's geometrical interpretation of $\epsilon(b)$ is the special case of the outer
surface, and the same interpretation applies to the CMB and to any other internal boundary,
each of which must be a level surface in a hydrostatic planet. This is why the profile of
$1/\epsilon$ plotted in the lecture, and computed for the two-layer model in Example~4, can be
read directly as the flattening of the Earth's internal structure: a two-layer Earth in
hydrostatic equilibrium would have a CMB with equatorial and polar radii differing by
$c/374\approx9\,\mathrm{km}$.
\end{solution}

\end{document}
